\textbf{Phần 1. Trắc nghiệm lựa chọn}
\setcounterpageref{ex}{0}

\begin{ex} %[3D1V08]
Một số có ba chữ số có chữ số hàng chục là $5$, chữ số hàng trăm gấp đôi chữ số hàng đơn vị. Biết tổng ba chữ số bằng $11$. Số đó là
\choice
{$254$}
{\True $452$}
{$653$}
{$854$}
\loigiai{
Tổng của chữ số hàng trăm và chữ số hàng đơn vị là:
$$11 - 5 = 6$$
Vì chữ số hàng trăm gấp đôi chữ số hàng đơn vị nên ta có sơ đồ:\\
- Chữ số hàng đơn vị: 1 phần\\
- Chữ số hàng trăm: 2 phần\\
Chữ số hàng đơn vị là:
$$6 : (2 + 1) = 2$$
Chữ số hàng trăm là:
$$2 \times 2 = 4$$
Vậy số cần tìm là $452$.\\
\textbf{Đáp án của câu hỏi là B}
}
\end{ex}

\begin{ex} %[3D1V08]
Tính nhanh giá trị của biểu thức: $B = (150 + 280 + 350) - (80 + 50)$.
\choice
{$600$}
{$620$}
{$680$}
{\True $650$}
\loigiai{
Nhóm các số để trừ tròn trăm:\\
$$B = (150 + 350 - 50) + (280 - 80) = 450 + 200 = 650$$
\textbf{Đáp án của câu hỏi là D}
}
\end{ex}

\begin{ex} %[3D1V08]
Tìm giá trị của $x$, biết: $5 \times x - 18 = 16 + 6$.
\choice
{\True $8$}
{$7$}
{$9$}
{$6$}
\loigiai{
Tính giá trị của vế phải trước:\\
$$5 \times x - 18 = 22$$
$$5 \times x = 22 + 18$$
$$5 \times x = 40$$
$$x = 40 : 5 = 8$$
\textbf{Đáp án của câu hỏi là A}
}
\end{ex}

\begin{ex} %[3D1V08]
Đoạn thẳng thứ nhất dài $45\text{ cm}$, đoạn thẳng thứ hai ngắn hơn đoạn thẳng thứ nhất $15\text{ cm}$. Đoạn thẳng thứ ba dài gấp đôi đoạn thẳng thứ hai. Độ dài đoạn thẳng thứ ba là
\choice
{$50\text{ cm}$}
{$70\text{ cm}$}
{\True $60\text{ cm}$}
{$40\text{ cm}$}
\loigiai{
Độ dài đoạn thẳng thứ hai là:
$$45 - 15 = 30\text{ (cm)}$$
Độ dài đoạn thẳng thứ ba là:
$$30 \times 2 = 60\text{ (cm)}$$
\textbf{Đáp án của câu hỏi là C}
}
\end{ex}

\begin{ex} %[3D1V08]
Cho bốn chữ số $0; 2; 5; 7$. Có thể viết được tất cả bao nhiêu số có 3 chữ số khác nhau và chia hết cho $5$ từ các chữ số đã cho?
\choice
{$8$ số}
{\True $10$ số}
{$12$ số}
{$6$ số}
\loigiai{
Một số chia hết cho 5 khi chữ số hàng đơn vị là 0 hoặc 5.\\
- \textbf{Trường hợp 1:} Chữ số hàng đơn vị là 0.\\
  + Hàng trăm có 3 cách chọn (từ 2; 5; 7).\\
  + Hàng chục có 2 cách chọn.\\
  Số lượng số: $3 \times 2 = 6$ (số).\\
- \textbf{Trường hợp 2:} Chữ số hàng đơn vị là 5.\\
  + Hàng trăm có 2 cách chọn (chọn 2 hoặc 7, vì khác 0).\\
  + Hàng chục có 2 cách chọn (bao gồm cả chữ số 0).\\
  Số lượng số: $2 \times 2 = 4$ (số).\\
Tổng cộng có: $6 + 4 = 10$ số thỏa mãn.\\
\textbf{Đáp án của câu hỏi là B}
}
\end{ex}

\begin{ex} %[3D1V08]
Cho dãy số cách đều: $1; 4; 7; 10; 13; \dots$ Số hạng thứ 10 của dãy số trên là
\choice
{$25$}
{$30$}
{$31$}
{\True $28$}
\loigiai{
Dãy số có khoảng cách giữa hai số liền nhau là: $4 - 1 = 3$ đơn vị.\\
Số hạng thứ 10 của dãy số là:
$$1 + (10 - 1) \times 3 = 1 + 27 = 28$$
\textbf{Đáp án của câu hỏi là D}
}
\end{ex}

\begin{ex} %[3D1V08]
Một hình tam giác có độ dài hai cạnh lần lượt là $16\text{ cm}$ và $24\text{ cm}$. Biết chu vi của hình tam giác đó bằng $60\text{ cm}$. Độ dài cạnh thứ ba của hình tam giác là
\choice
{\True $20\text{ cm}$}
{$18\text{ cm}$}
{$22\text{ cm}$}
{$25\text{ cm}$}
\loigiai{
Tổng độ dài hai cạnh đã biết là:
$$16 + 24 = 40\text{ (cm)}$$
Độ dài cạnh thứ ba của hình tam giác là:
$$60 - 40 = 20\text{ (cm)}$$
\textbf{Đáp án của câu hỏi là A}
}
\end{ex}

\begin{ex} %[3D1V08]
Trong một phép trừ có số bị trừ bằng $450$. Nếu giữ nguyên số bị trừ và bớt số trừ đi $25$ đơn vị thì được hiệu mới bằng $185$. Số trừ ban đầu là
\choice
{$280$}
{$300$}
{\True $290$}
{$265$}
\loigiai{
Khi bớt ở số trừ đi 25 đơn vị thì hiệu tăng thêm 25 đơn vị.\\
Hiệu ban đầu của phép trừ là:
$$185 - 25 = 160$$
Số trừ ban đầu là:
$$450 - 160 = 290$$
\textbf{Đáp án của câu hỏi là C}
}
\end{ex}

\begin{ex} %[3D1V08]
Mai mua $4$ chiếc bút chì, mỗi chiếc có giá $5$ nghìn đồng và mua thêm $1$ quyển vở giá $8$ nghìn đồng. Mai đưa cho cô bán hàng tờ tiền $50$ nghìn đồng. Hỏi cô bán hàng phải trả lại Mai bao nhiêu nghìn đồng?
\choice
{$20$ nghìn đồng}
{\True $22$ nghìn đồng}
{$24$ nghìn đồng}
{$18$ nghìn đồng}
\loigiai{
Số tiền Mai mua 4 chiếc bút chì là:
$$4 \times 5 = 20\text{ (nghìn đồng)}$$
Tổng số tiền Mai mua đồ là:
$$20 + 8 = 28\text{ (nghìn đồng)}$$
Số tiền cô bán hàng trả lại Mai là:
$$50 - 28 = 22\text{ (nghìn đồng)}$$
\textbf{Đáp án của câu hỏi là B}
}
\end{ex}

\begin{ex} %[3D1V08]
Hai số có tổng bằng $450$ và hiệu của chúng bằng $50$. Số lớn là
\choice
{$200$}
{$300$}
{$225$}
{\True $250$}
\loigiai{
Số lớn là:
$$(450 + 50) : 2 = 250$$
\textbf{Đáp án của câu hỏi là D}
}
\end{ex}

\begin{ex} %[3D1V08]
Một đoạn đường dài $32\text{ m}$, người ta trồng các cây cách đều nhau $4\text{ m}$, ở cả hai đầu đường đều có trồng cây. Hỏi đoạn đường đó trồng tất cả bao nhiêu cây?
\choice
{\True $9$ cây}
{$8$ cây}
{$10$ cây}
{$7$ cây}
\loigiai{
Số khoảng cách giữa các cây là:
$$32 : 4 = 8\text{ (khoảng)}$$
Vì cả hai đầu đường đều có cây nên số cây trồng được là:
$$8 + 1 = 9\text{ (cây)}$$
\textbf{Đáp án của câu hỏi là A}
}
\end{ex}

\begin{ex} %[3D1V08]
Một số tự nhiên khi chia cho $4$ thì được thương bằng $9$ và số dư là số dư lớn nhất có thể có. Khi lấy số đó chia cho $5$ thì số dư là
\choice
{$2$}
{$3$}
{\True $4$}
{$1$}
\loigiai{
Trong phép chia cho 4, số dư lớn nhất có thể là $3$.\\
Số tự nhiên đó là:
$$4 \times 9 + 3 = 36 + 3 = 39$$
Lấy 39 chia cho 5 ta có:
$$39 = 5 \times 7 + 4$$
Vậy khi chia cho 5 thì số dư là $4$.\\
\textbf{Đáp án của câu hỏi là C}
}
\end{ex}

\textbf{Phần 2. Trắc nghiệm đúng sai}
\setcounterpageref{ex}{0}

\begin{ex} %[3D1V08]
Xét tính đúng sai của các nhận định sau:
\choiceTF
{\True Trong một phép chia cho $5$, số dư lớn nhất có thể có bằng $4$.}
{\True Chu vi của hình tam giác có ba cạnh bằng nhau mỗi cạnh $6\text{ cm}$ bằng chu vi của hình chữ nhật có chiều dài $5\text{ cm}$ và chiều rộng $4\text{ cm}$.}
{Khi cộng một số tự nhiên bất kì với $0$ thì luôn được kết quả bằng $0$.}
{\True Tích của số $2$ với bất kì một số lẻ nào cũng luôn là một số chẵn.}
\loigiai{
\begin{itemize}
    \item Ý a) Đúng: Số dư luôn nhỏ hơn số chia nên số dư lớn nhất là 4.
    \item Ý b) Đúng: Chu vi tam giác: $6 \times 3 = 18\text{ (cm)}$. Chu vi hình chữ nhật: $(5 + 4) \times 2 = 18\text{ (cm)}$.
    \item Ý c) Sai: Khi cộng một số với 0 thì bằng chính số đó ($a + 0 = a$).
    \item Ý d) Đúng: Số 2 nhân với bất kì số nào cũng cho kết quả là số chẵn.
\end{itemize}
\begin{itemize}
    \item a) \textbf{Đúng}.
    \item b) \textbf{Đúng}.
    \item c) \textbf{Sai}.
    \item d) \textbf{Đúng}.
\end{itemize}
}
\end{ex}

\begin{ex} %[3D1V08]
Một đội công nhân phải sửa đoạn đường dài $300\text{ m}$ trong ba ngày. Ngày thứ nhất sửa được $80\text{ m}$, ngày thứ hai sửa được số mét đường gấp đôi ngày thứ nhất. Xét tính đúng sai:
\choiceTF
{\True Ngày thứ hai đội công nhân sửa được $160\text{ m}$ đường.}
{\True Sau hai ngày đầu, đội công nhân đã sửa được tất cả $240\text{ m}$ đường.}
{Ngày thứ ba đội công nhân phải sửa $70\text{ m}$ đường nữa mới hoàn thành công việc.}
{\True Ngày thứ hai đội công nhân sửa được nhiều hơn ngày thứ ba là $100\text{ m}$ đường.}
\loigiai{
\begin{itemize}
    \item Ý a) Đúng: Ngày thứ hai sửa được: $80 \times 2 = 160\text{ (m)}$.
    \item Ý b) Đúng: Sau hai ngày đầu sửa được: $80 + 160 = 240\text{ (m)}$.
    \item Ý c) Sai: Ngày thứ ba cần sửa là: $300 - 240 = 60\text{ (m)}$ (không phải 70 m).
    \item Ý d) Đúng: Ngày thứ hai sửa nhiều hơn ngày thứ ba là: $160 - 60 = 100\text{ (m)}$.
\end{itemize}
\begin{itemize}
    \item a) \textbf{Đúng}.
    \item b) \textbf{Đúng}.
    \item c) \textbf{Sai}.
    \item d) \textbf{Đúng}.
\end{itemize}
}
\end{ex}

\textbf{Phần 3. Trả lời ngắn}
\setcounterpageref{ex}{0}

\begin{ex} %[3D1V08]
Tính nhanh giá trị của biểu thức: $135 + 246 + 165 + 154$.
\shortans{700}
\loigiai{
\textbf{Đáp số:} $700$\\
Nhóm các số tròn trăm:\\
$$(135 + 165) + (246 + 154) = 300 + 400 = 700$$
}
\end{ex}

\begin{ex} %[3D1V08]
Tìm giá trị của $x$, biết: $x : 3 + 45 = 54$.
\shortans{27}
\loigiai{
\textbf{Đáp số:} $27$\\
Xem $x : 3$ là số hạng chưa biết:\\
$$x : 3 = 54 - 45$$
$$x : 3 = 9$$
$$x = 9 \times 3 = 27$$
}
\end{ex}

\begin{ex} %[3D1V08]
Một khúc gỗ dài $20\text{ m}$ được người ta cưa thành các đoạn ngắn dài $4\text{ m}$. Hỏi người thợ phải cưa tất cả bao nhiêu lần?
\shortans{4}
\loigiai{
\textbf{Đáp số:} $4$\\
Số đoạn gỗ cưa được là: $20 : 4 = 5$ (đoạn).\\
Số lần cưa cần thực hiện là: $5 - 1 = 4$ (lần cưa).
}
\end{ex}

\begin{ex} %[3D1V08]
Hai thùng đựng tất cả $160\text{ lít}$ nước mắm. Thùng thứ nhất đựng nhiều hơn thùng thứ hai $20\text{ lít}$. Hỏi thùng thứ nhất đựng bao nhiêu lít nước mắm?
\shortans{90}
\loigiai{
\textbf{Đáp số:} $90$\\
Thùng thứ nhất đựng số lít nước mắm là:
$$(160 + 20) : 2 = 90\text{ (lít)}$$
}
\end{ex}

\begin{ex} %[3D1V08]
Cho ba chữ số $2; 4; 7$. Tìm hiệu giữa số lớn nhất có 3 chữ số khác nhau và số bé nhất có 3 chữ số khác nhau được lập từ ba chữ số đã cho.
\shortans{495}
\loigiai{
\textbf{Đáp số:} $495$\\
- Số lớn nhất có 3 chữ số khác nhau là: $742$.\\
- Số bé nhất có 3 chữ số khác nhau là: $247$.\\
Hiệu giữa hai số là:
$$742 - 247 = 495$$
}
\end{ex}

\begin{ex} %[3D1V08]
Dãy số chẵn liên tiếp $2; 4; 6; 8; \dots; 40$ có tất cả bao nhiêu số hạng?
\shortans{20}
\loigiai{
\textbf{Đáp số:} $20$\\
Số số hạng của dãy là:
$$(40 - 2) : 2 + 1 = 38 : 2 + 1 = 20\text{ (số hạng)}$$
}
\end{ex}
